% =====================================================================
%  refs.tex —— 参考文献（原文 pp.50–54）。条目按底本原样保留英文。
% =====================================================================
\addcontentsline{toc}{section}{参考文献}
\begin{thebibliography}{99}
\bibitem{1} E. Witten, “Geometric Langlands from six dimensions,” arXiv:0905.2720 [hep-th].
\bibitem{2} D. Gaiotto, “N = 2 dualities,” arXiv:0904.2715 [hep-th].
\bibitem{3} D. Gaiotto, G. W. Moore, and A. Neitzke, “Wall-crossing, hitchin Systems, and the WKB approximation,” arXiv:0907.3987 [hep-th].
\bibitem{4} P. C. Argyres and N. Seiberg, “S-Duality in N = 2 supersymmetric gauge theories,” JHEP 12 (2007) 088, arXiv:0711.0054 [hep-th].
\bibitem{5} J. A. Minahan and D. Nemeschansky, “An N = 2 superconformal fixed point with E6 global symmetry,” Nucl. Phys. B482 (1996) 142–152, arXiv:hep-th/9608047.
\bibitem{6} P. C. Argyres and M. R. Douglas, “New phenomena in SU(3) supersymmetric gauge theory,” Nucl.Phys. B448 (1995) 93–126, arXiv:hep-th/9505062 [hep-th].
\bibitem{7} P. C. Argyres, M. Plesser, and A. D. Shapere, “The Coulomb phase of N=2 supersymmetric QCD,” Phys.Rev.Lett. 75 (1995) 1699–1702, arXiv:hep-th/9505100 [hep-th].
\bibitem{8} N. Hitchin, “The self-duality equation on a riemann surface,” Proc.London Math.Soc. (3)55 (1987) 59–126.
\bibitem{9} N. Hitchin, “Stable bundles and integrable system,” Duke Math. J. (1)54 (1987) 91–114.
\bibitem{10} O. Biquard and P. Boalch, “Wild non-abelian Hodge theory on curves,” Compos. Math. 140 (2004) no. 1, 179–204. http://dx.doi.org/10.1112/S0010437X03000010.
\bibitem{11} P. C. Argyres, M. Ronen Plesser, N. Seiberg, and E. Witten, “New N=2 superconformal field theories in four-dimensions,” Nucl.Phys. B461 (1996) 71–84, arXiv:hep-th/9511154 [hep-th].
\bibitem{12} T. Eguchi, K. Hori, K. Ito, and S.-K. Yang, “Study of N = 2 Superconformal Field Theories in 4 Dimensions,” Nucl. Phys. B471 (1996) 430–444, arXiv:hep-th/9603002.
\bibitem{13} T. Eguchi and K. Hori, “N = 2 superconformal field theories in 4 dimensions and A- D-E classification,” arXiv:hep-th/9607125.
\bibitem{14} T. Banks, M. R. Douglas, and N. Seiberg, “Probing F-theory with branes,” Phys. Lett. B387 (1996) 278–281, arXiv:hep-th/9605199.
\bibitem{15} M. R. Douglas, D. A. Lowe, and J. H. Schwarz, “Probing F-theory with multiple branes,” Phys. Lett. B394 (1997) 297–301, arXiv:hep-th/9612062.
\bibitem{16} S. Gukov and A. Kapustin, “New N = 2 superconformal field theories from M/F theory orbifolds,” Nucl. Phys. B545 (1999) 283–308, arXiv:hep-th/9808175.
\bibitem{17} S. Cecotti, A. Neitzke, and C. Vafa, “R-twisting and 4d/2d correspondences,” arXiv:1006.3435 [hep-th].
\bibitem{18} J. Seo and K. Dasgupta, “Argyres-Douglas Loci, Singularity Structures and Wall- Crossings in Pure N=2 Gauge Theories with Classical Gauge Groups,” arXiv:1203.6357 [hep-th].
\bibitem{19} S. M. Kuzenko and S. Theisen, “Correlation functions of conserved currents in N = 2 superconformal theory,” Class. Quant. Grav. 17 (2000) 665–696, arXiv:hep-th/9907107.
\bibitem{20} A. D. Shapere and Y. Tachikawa, “Central charges of N=2 superconformal field theories in four dimensions,” JHEP 0809 (2008) 109, arXiv:0804.1957 [hep-th].
\bibitem{21} O. Aharony and Y. Tachikawa, “A holographic computation of the central charges of d=4, N=2 SCFTs,” JHEP 01 (2008) 037, arXiv:0711.4532 [hep-th].
\bibitem{22} D. Gaiotto, N. Seiberg, and Y. Tachikawa, “Comments on scaling limits of 4d N=2 theories,” JHEP 01 (2011) 078, arXiv:1011.4568 [hep-th].
\bibitem{23} D. Nanopoulos and D. Xie, “Hitchin equation, singularity, and N = 2 superconformal field theories,” JHEP 03 (2010) 043, arXiv:0911.1990 [hep-th].
\bibitem{24} A. Kapustin, “Solution of N = 2 gauge theories via compactification to three dimensions,” Nucl. Phys. B534 (1998) 531–545, arXiv:hep-th/9804069.
\bibitem{25} A. Kapustin and E. Witten, “Electric-Magnetic duality and the Geometric Langlands Program,” arXiv:hep-th/0604151.
\bibitem{26} E. Witten, “Gauge theory and wild ramification,” arXiv:0710.0631 [hep-th].
\bibitem{27} W. Wasow, Asymptotic expansions for ordinary differential equations. Dover Publications Inc., New York, 1987. Reprint of the 1976 edition.
\bibitem{28} P. P. Boalch, “G-bundles, isomonodromy, and quantum Weyl groups,” Int. Math. Res. Not. (2002) no. 22, 1129–1166. http://dx.doi.org/10.1155/S1073792802111081.
\bibitem{29} S. Gukov and E. Witten, “Gauge Theory, ramification, and the Geometric Langlands Program,” arXiv:hep-th/0612073.
\bibitem{30} D. Xie, “Network, cluster coordinates and N = 2 theory II: Irregular singularity,” arXiv:1207.6112 [hep-th].
\bibitem{31} S. Cecotti and C. Vafa, “Classification of complete N=2 supersymmetric theories in 4 dimensions,” arXiv:1103.5832 [hep-th].
\bibitem{32} D. Nanopoulos and D. Xie, “Hitchin equation, irregular Singularity, and N = 2 asymptotical free theories,” arXiv:1005.1350 [hep-th].
\bibitem{33} A. D. Shapere and C. Vafa, “BPS structure of Argyres-Douglas superconformal theories,” arXiv:hep-th/9910182.
\bibitem{34} M. Alim, S. Cecotti, C. Cordova, S. Espahbodi, A. Rastogi, and C. Vafa, “BPS Quivers and Spectra of Complete N=2 Quantum Field Theories,” arXiv:1109.4941 [hep-th].
\bibitem{35} M. Alim, S. Cecotti, C. Cordova, S. Espahbodi, A. Rastogi, and C. Vafa, “N=2 Quantum Field Theories and Their BPS Quivers,” arXiv:1112.3984 [hep-th].
\bibitem{36} D. Gaiotto, G. W. Moore, and A. Neitzke, “Framed BPS states,” arXiv:1006.0146 [hep-th].
\bibitem{37} D. Gaiotto, G. W. Moore, and A. Neitzke, “Wall-Crossing in Coupled 2d-4d Systems,” arXiv:1103.2598 [hep-th].
\bibitem{38} F. Benini, Y. Tachikawa, and D. Xie, “Mirrors of 3d Sicilian theories,” JHEP 1009 (2010) 063, arXiv:1007.0992 [hep-th].
\bibitem{39} P. Boalch, “Irregular connections and Kac-Moody root systems,” arXiv:0806.1050.
\bibitem{40} D. Nanopoulos and D. Xie, “More three dimensional mirror pairs,” arXiv:1011.1911 [hep-th].
\bibitem{41} S. A. Cherkis and A. Kapustin, “Periodic monopoles with singularities and N = 2 super-QCD,” Commun. Math. Phys. 234 (2003) 1–35, arXiv:hep-th/0011081.
\bibitem{42} D. Gaiotto, “unpublished,”.
\bibitem{43} S. Fomin, M. Shapiro, and D. Thurston, “Cluster algebras and triangulated surfaces. part i: Cluster complexes,” ACTA MATHEMATICA 201 (2008) 83. arXiv.org:math/0608367.
\bibitem{44} K. A. Intriligator and N. Seiberg, “Mirror symmetry in three dimensional gauge theories,” Phys. Lett. B387 (1996) 513–519, arXiv:hep-th/9607207.
\bibitem{45} E. Witten, “Solutions of four-dimensional field theories via M-theory,” Nucl. Phys. B500 (1997) 3–42, arXiv:hep-th/9703166.
\bibitem{46} F. Benini, S. Benvenuti, and Y. Tachikawa, “Webs of five-Branes and N = 2 superconformal field theories,” JHEP 09 (2009) 052, arXiv:0906.0359 [hep-th].
\bibitem{47} S. Cecotti and M. Del Zotto, “On Arnold’s 14 ‘exceptional’ N=2 superconformal gauge theories,” JHEP 10 (2011) 099, arXiv:1107.5747 [hep-th].
\bibitem{48} M. Del Zotto, “More Arnold’s N = 2 superconformal gauge theories,” JHEP 11 (2011) 115, arXiv:1110.3826 [hep-th].
\bibitem{49} S. Katz, P. Mayr, and C. Vafa, “Mirror symmetry and exact solution of 4D N = 2 gauge theories. I,” Adv. Theor. Math. Phys. 1 (1998) 53–114, arXiv:hep-th/9706110.
\bibitem{50} H. Lin, O. Lunin, and J. M. Maldacena, “Bubbling AdS space and 1/2 BPS geometries,” JHEP 10 (2004) 025, arXiv:hep-th/0409174.
\bibitem{51} D. Gaiotto and J. Maldacena, “The gravity duals of N = 2 superconformal field theories,” arXiv:0904.4466 [hep-th].
\bibitem{52} O. Aharony, A. Fayyazuddin, and J. M. Maldacena, “The large N limit of N = 2,1 field theories from three- branes in F-theory,” JHEP 07 (1998) 013, arXiv:hep-th/9806159.
\bibitem{53} D. Xie, “Network, Cluster coordinates and N=2 theory I,” arXiv:1203.4573 [hep-th].
\bibitem{54} L. F. Alday, D. Gaiotto, and Y. Tachikawa, “Liouville correlation functions from four-dimensional gauge theories,” Lett. Math. Phys. 91 (2010) 167–197, arXiv:0906.3219 [hep-th].
\bibitem{55} D. Gaiotto and J. Teschner, “Irregular singularities in Liouville theory,” arXiv:1203.1052 [hep-th].
\bibitem{56} G. Bonelli, K. Maruyoshi, and A. Tanzini, “Wild Quiver Gauge Theories,” JHEP 02 (2012) 031, arXiv:1112.1691 [hep-th].
\bibitem{57} N. Nekrasov and E. Witten, “The Omega deformation, branes, integrability, and liouville Theory,” JHEP 09 (2010) 092, arXiv:1002.0888 [hep-th].
\bibitem{58} D. Gaiotto, “Asymptotically free N=2 theories and irregular conformal blocks,” arXiv:0908.0307 [hep-th].
\bibitem{59} N. Nekrasov, A. Rosly, and S. Shatashvili, “Darboux coordinates, Yang-Yang functional, and gauge theory,” Nucl. Phys. Proc. Suppl. 216 (2011) 69–93, arXiv:1103.3919 [hep-th].
\bibitem{60} V. V. Fock and A. B. Goncharov, “The quantum dilogarithm and representations quantum cluster varieties,” arXiv:0702397 [math].
\bibitem{61} B. Feigin, E. Frenkel, and V. Toledano-Laredo, “Gaudin models with irregular singularities,” 2006. doi:10.1016/j.aim.2009.09.007.
\bibitem{62} P. Boalch, “Simply-laced isomonodromy systems,” 2011. http://www.citebase.org/abstract?id=oai:arXiv.org:1107.0874.
\bibitem{63} S. Cecotti, C. Cordova, and C. Vafa, “Braids, Walls, and Mirrors,” arXiv:1110.2115 [hep-th].
\bibitem{64} Z. Komargodski and A. Schwimmer, “On Renormalization Group Flows in Four Dimensions,” JHEP 12 (2011) 099, arXiv:1107.3987 [hep-th].
\bibitem{65} N. Arkani-Hamed, J. L. Bourjaily, F. Cachazo, A. B. Goncharov, A. Postnikov, and J. Trnka, Scattering Amplitudes and the Positive Grassmannian. Cambridge University Press, 2012. arXiv:1212.5605 [hep-th]. http://inspirehep.net/record/1208741/files/arXiv:1212.5605.pdf.
\bibitem{66} A. Postnikov, “Total positivity, grassmannians, and networks,”. arXiv.org:math/0609764.
\bibitem{67} A. Gadde, L. Rastelli, S. S. Razamat, and W. Yan, “Gauge Theories and Macdonald Polynomials,” arXiv:1110.3740 [hep-th].
\end{thebibliography}
